\section{Discussion}
\label{sec:disc}

\paragraph{Whether to build an internal judge is a question about the model.}
The result we did not expect is that the answer changes by family, and changes completely, with the
separation explained neither by labelled failures, size, attention design nor release date. What follows is a
measurement rather than a rule of thumb: score confidence on a few hundred labelled calls from the model to be
deployed, and build an internal judge only if that score is poor. Two earlier conclusions of this study,
drawn when it contained one family, were overturned by more: that small models cannot abstain when no tool
applies, and that confidence is uniformly useless (Appendix~\ref{app:arch}).

\paragraph{Resolution, not the readout.}
Within the attention tier the head-averaged readouts fail and the per-head readouts succeed on identical
data, width-matched controls attribute the difference to resolution, and the strongest published member of
the tier reads an in-degree sequence, so the two best attention detectors agree on almost nothing except
keeping heads apart. Within the residual tier the detectors that read designated positions lead and those
that pool over the span trail. In both tiers the detectors that preserve locality win.

\paragraph{Negative results worth stating.}
We do not beat the attention-only state of the art; at matched protocol the two are within their paired
intervals, and our contribution to that tier is an explanation and a compression. The two families are not
redundant, yet nothing we tried converts that into a better detector: use the judge your access allows and
spend the labelling budget on its threshold. The largest effects we met while building the study were defects
in the measurement chain, each producing a plausible wrong number, and the confound panel and the
failure-mode stratification were the only reason we noticed (Appendix~\ref{app:validity}).

\paragraph{Limitations.}
Models run from 0.8B to 4B parameters, chosen so a dozen runs with per-head eigendecompositions fit one
consumer GPU; how the family split behaves at scale is the first experiment we would add. Detectors are
trained per model. The multi-turn evaluation is teacher-forced and cannot show what a wrong tool result
does. Labels compare a call to a benchmark's ground truth, a proxy for harm. The family split is confounded
with the call format the families emit. Every attention feature is computed on a symmetrised graph and cannot
see the direction of attention flow \citep{dahlem2026transport}; the antisymmetric part and the sink
statistics of \citet{binkowski2026sinks} are unread channels. The attention tier's cost was measured on one
GPU, not a production stack.
